\begin{solapartat}
\punts{0,25} La longitud d'ona llindar $\lambda_0 = 650\cdot 10^{-9}\ \mathrm{m}$ és la longitud d'ona més alta per la qual tenim efecte fotoelèctric. Correspon a una freqüència: $f_0 = \frac{c}{\lambda} = \frac{3,0\cdot 10^{8}}{650\cdot 10^{-9}} = 4,62\times 10^{14}\ \mathrm{Hz}$.

\punts{0,25} El treball d'extracció del metall és l'energia mínima per extreure l'electró: $hf = E_C + W_0$, per tant, per $E_c = 0$, $W_0 = hf_0$ i fent el càlcul obtenim:
\[ W_0 = hf_0 = 6,626\times 10^{-34}\cdot 4,62\times 10^{14} = 3,06\cdot 10^{-19}\ \mathrm{J} = 1,91\ \mathrm{eV} \]

\punts{0,25} El potencial de frenada és el voltatge que necessitem aplicar per frenar els electrons amb més energia cinètica: $\Delta U = |e|\Delta V = Ec$

% A l'original: "Hzi. I," (probablement "Hz. I,").
\punts{0,25} Calculem l'energia cinètica amb què surten els electrons per 300~nm. Primer, calculem la freqüència corresponent: $f = \frac{c}{\lambda} = \frac{3,0\cdot 10^{8}}{300\cdot 10^{-9}} = 1,00\times 10^{15}\ \mathrm{Hz}$. I, seguidament, l'energia cinètica amb què surten és:
\[ E_C = hf - W_0 = 6,626\times 10^{-34}\cdot 1,00\times 10^{15} - 3,06\cdot 10^{-19} = 3,57\cdot 10^{-19}\ \mathrm{J} \]

\punts{0,25} L'energia cinètica es contraresta amb una energia potencial elèctrica d'igual mòdul: $\Delta U = q\cdot\Delta V$ i, per tant, el potencial de frenada ha de ser $\Delta V = \frac{Ec}{|e|} = \frac{3,57\cdot 10^{-19}}{1,602\cdot 10^{-19}} = 2,23\ \mathrm{V}$.
\end{solapartat}

\begin{solapartat}
\punts{0,5} A partir de l'equació de l'energia cinètica ja utilitzada tenim $E_C = \frac{1}{2}mv^2 = hf - W_0$, i posant la freqüència en funció de la longitud d'ona s'obté: $v = \sqrt{\frac{2}{m}\left(\frac{hc}{\lambda} - W_0\right)}$.

Introduint els valors corresponents:
\[ v(\lambda) = \sqrt{\frac{2}{9,11\cdot 10^{-31}}\left(\frac{6,626\cdot 10^{-34}\cdot 3\cdot 10^{8}}{\lambda} - 3,06\cdot 10^{-19}\right)}\ \mathrm{\frac{m}{s}}, \]
per tant, l'expressió és $v(\lambda) = \sqrt{\frac{4,36\cdot 10^{5}}{\lambda} - 6,72\cdot 10^{11}}\ \mathrm{\frac{m}{s}}$

\punts{0,25} La velocitat per 500~nm: \quad $v(500\ \mathrm{nm}) = \sqrt{\frac{4,36\cdot 10^{5}}{500\cdot 10^{-9}} - 6,72\cdot 10^{11}} = 4,47\cdot 10^{5}\ \mathrm{\frac{m}{s}}$

\punts{0,5} La longitud d'ona de De Broglie la calculem a partir de la quantitat de moviment de l'electró i la relació de De Broglie: \ $p = mv = \frac{h}{\lambda}$.

% DUBTE: 1,63·10^-9 m = 1,63 nm, no 1,67 nm (errada de la pauta, transcrita tal com és).
Per tant, $\lambda = \frac{h}{mv} = \frac{6,626\cdot 10^{-34}}{9,11\cdot 10^{-31}\cdot 4,47\cdot 10^{5}} = 1,63\cdot 10^{-9}\ \mathrm{m} = 1,67\ \mathrm{nm}$.
\end{solapartat}
